\documentclass[10pt,a4paper]{amsart}
\usepackage[margin=19mm]{geometry}
\usepackage{amsmath,amssymb,mathtools}
\usepackage[scr=rsfs,cal=euler]{mathalfa}
\usepackage{xfrac}
\usepackage[unicode,hidelinks,bookmarks=false]{hyperref}
\newcommand{\ie}{i.e.\ }
\newcommand{\cf}{cf.\ }
\newcommand{\resp}{respectively\ }
\newcommand{\Subsection}{\subsection}
\providecommand{\nobreakdash}{\nobreak\hbox{-}}
\setlength{\emergencystretch}{2em}
\multlinegap0pt
\usepackage{amssymb}
\usepackage{xcolor}
\usepackage{bm}
\usepackage{dsfont}
\usepackage{float}
\usepackage{xfrac}
\usepackage{tikz}
\newtheorem{lem}{Lemma}
\title{Fourier Dimension in Inhomogeneous Duffin--Schaeffer Conjecture}
\author{Bo Tan, Qing-Long Zhou}
\date{Source: arXiv:2604.13868v2 (CC BY 4.0)}
\begin{document}
\maketitle

\noindent\emph{Source and licence.} Fourier Dimension in Inhomogeneous Duffin--Schaeffer Conjecture. Version arXiv:2604.13868v2 (CC BY 4.0), distributed by arXiv under the Creative Commons Attribution 4.0 International licence, http://creativecommons.org/licenses/by/4.0/, as recorded in the arXiv licence metadata for this exact version and shown on the abstract page https://arxiv.org/abs/2604.13868v2. Full licence notice: \texttt{licenses/arxiv-2604.13868-CC-BY-4.0.txt}.\par\smallskip

\noindent\textit{Editorial scope.} Counted unit: the article's lem gM\_Fourier (source0.tex lines 727--742). The article's complete original proof of it is delivered with it (source0.tex lines 744--790, one \texttt{proof} environment of 47 lines). No mathematics is written, repaired or re-derived: the statement and the proof are the article's own lines, in the article's own notation, and no retained line is substituted or re-flowed.  Retained uncounted as the source-local context the proof needs: lines 490--492; lines 500--504; lines 706--708.  Nothing was omitted from any retained span, so the package carries no omission record.  No retained line contains a citation, so the wrapper carries no bibliography and no bibliography entry.  No retained line contains a figure environment, an image-inclusion command or drawing code, so no image is delivered and the package declares no asset dependency.  Editorial formatting only, in addition: an AMS \texttt{amsart} preamble that restates the article's own declarations and macros used by the retained lines, namely the environments \texttt{lem} on separate counters, together with the standard packages those lines need.  The retained environments print in this document as Lemma 1 in order of appearance, whereas the article prints the same items with the numbering of its own full document.  The complete native source of the article as distributed in the arXiv e-print for this exact version is bundled unchanged as \texttt{sources/198504/source0.tex}.
\begin{equation}\label{cardIq}
  \# \mathcal{I}_q = \phi(q) \prod_{p \mid (q, B_q)} \left(1 + \frac{1}{p-1}\right) = q \prod_{ {p \mid q,~p \nmid B_q}} \bigl(1 - p^{-1}\bigr),
\end{equation}

\begin{lem}\label{grs}
Let $A, B$ be coprime integers with $B>0.$ Then for $q,k\in\mathbb{N},$ we have
\[\left|\sum_{\substack{0\le p\le q-1 \\ \gcd(Bp+A,q)=1}}e(-k\cdot \frac{p}{q})\right|\ll \gcd(q,k)\cdot \log q,\]
where the implied constant in Vinogradov’s notation is independent of $A,B,q,k.$
\end{lem}

\begin{equation}\label{Q(M)}
\# Q(M) \;\gg\; \frac{M}{(\log M)^{2}}.
\end{equation}

\begin{lem}\label{gM_Fourier}
For $\varepsilon>0$, there exists $M_0^{\ast}>0$ such that for all $M\in\mathcal{M}$ with $M\ge M_0^{\ast}$, the Fourier coefficients of $g_M$ satisfy the following estimates:
\begin{enumerate}
\item[(1)] $\widehat{g_M}(0)=1$.
\item[(2)] For $l\in \mathbb{Z}\setminus \{0\}$,
\[
 \widehat{g_M}(l)  \ll
\begin{cases}
\dfrac{(\log M)^{5}\,\tau(l)}{M} & \text{if } 1\le |l|\le 3 M^{1/\eta}, \\[10pt]
|l|^{-(\eta-\varepsilon)} & \text{if } |l| > M^{1/\eta}.
\end{cases}
\]
%And thus $\widehat{g_M}(l)  \ll M^{-1+}$

\end{enumerate}
\end{lem}

\begin{proof} A straightforward computation of the Fourier coefficient yields
\begin{align*}
\widehat{g_M}(l)
&= \int_{0}^{1} \frac{1}{\#Q(M)} \sum_{q\in Q(M)} \sum_{k\in\mathbb{Z}} \widehat{\Phi_{q,\theta}}(k) e(kx) \, e(-lx) \, \mathrm{d}x \\
&= \frac{1}{\#Q(M)} \sum_{q\in Q(M)} \widehat{\Phi_{q,\theta}}(l) \\
&= \frac{1}{\#Q(M)} \sum_{q\in Q(M)} \frac{1}{\#\mathcal{I}_q} \sum_{p\in\mathcal{I}_q} e\!\left(l\cdot\frac{p+\theta(q)}{q}\right) \widehat{f}\!\left(\frac{\psi(q)}{q}l\right).
\end{align*}

\noindent(1) When \(l = 0\), since \(\widehat{f}(0) = 1\), we have \(\widehat{g_M}(0) = 1\).

\smallskip

\noindent(2) When \(l \neq 0\), We treat the two cases in the lemma separately.

\medskip

\noindent{}~~\underline{\textbf{Case 1: $1\le |l|\le 3M^{1/\eta}$.}} We bound the Fourier coefficients in the following way.
\begin{align*}
\bigl|\widehat{g_M}(l)\bigr|
&\le \frac{(\log M)^2}{M} \sum_{q\in Q(M)} \frac{1}{q\prod_{p\le q}(1-p^{-1})} \cdot \Bigl|\sum_{p\in\mathcal{I}_q} e\!\bigl(l\cdot \tfrac{p}{q}\bigr)\Bigr| \\
&\ll \frac{(\log M)^2}{M} \sum_{q\in Q(M)} \frac{(\log q)^2 \gcd(q,l)}{q} \\
&\ll \frac{(\log M)^4}{M} \sum_{q=1}^{(2M)\!^{\frac{1}{\eta}}} \frac{\gcd(q,l)}{q} \\
&\le \frac{(\log M)^4}{M} \sum_{d\mid l} \sum_{\substack{1\le k\le (2M)\!^{\frac{1}{\eta}}/d \\ (k,l/d)=1}} \frac{d}{dk} \\
&\ll \frac{(\log M)^5}{M} \, \tau(l),
\end{align*}
where the first inequality is a consequence of the bound \(|\widehat{f}(\xi)| \le 1\) (\(\xi \in \mathbb{R}\)) and the counting estimates \eqref{cardIq} and \eqref{Q(M)};
the second inequality follows from Lemma \ref{grs} and Mertens' third theorem\footnote{Mertens' third theorem states that $\lim_{q\to\infty} \log q \prod_{p\le q}(1-p^{-1}) = e^{-\gamma}$, where $\gamma$ is the Euler-Mascheroni constant.}; the third inequality holds because
\[\frac{1}{2M}\le \left(\frac{\psi(q)}{q}\right)^{\eta}< \left(\frac{1}{q}\right)^{\eta} \qquad \implies \qquad \log q \ll_{\eta} \log M.\]

\medskip

\noindent{}~~\underline{\textbf{Case 2: $|l| > M^{1/\eta}$.}}
Since $|\widehat{f}(\xi)| \ll (1+|\xi|)^{-K}$ with $K>\eta$, we have
\[
\bigl|\widehat{g_M}(l)\bigr|
\le \frac{(\log M)^2}{M} \sum_{q\in Q(M)} \frac{1}{q\prod_{p\le q}(1-p^{-1})} \cdot \Bigl|\sum_{p\in\mathcal{I}_q} e\!\bigl(l\cdot \tfrac{p}{q}\bigr)\Bigr| \cdot \frac{1}{\bigl(1+\frac{\psi(q)}{q}|l|\bigr)^{\eta}}.
\]
Since $({\psi(q)}/{q})^{-\eta}\asymp M$, we deduce that  $(1+ |l|{\psi(q)}/{q})^{-\eta}\ll M|l|^{-\eta}$. Proceeding as in Case~1 and noting that \(|l| > M^{1/\eta}\), we have for sufficiently large \(M\),
\begin{align*}
\bigl|\widehat{g_M}(l)\bigr|
\ll \frac{(\log M)^5}{M} \, \tau(l) \cdot M \, |l|^{-\eta}
\le |l|^{-(\eta-\varepsilon)}.
\end{align*}
This completes the proof.


\end{proof}

\end{document}
